\documentclass{article}
\usepackage[T1]{fontenc}
\usepackage{lmodern}
\usepackage{graphicx}
\usepackage{amsmath,amssymb}
\usepackage[a4paper,margin=2cm]{geometry}
\usepackage[absolute,overlay]{textpos}
\setlength{\TPHorizModule}{1mm}
\setlength{\TPVertModule}{1mm}
\usepackage[indonesian]{babel}
\usepackage{hyperref}
\setlength{\emergencystretch}{2em}
% unit: Halaman 3

\begin{document}

% === Halaman 3 (llm) ===

\section*{Sejarah DES}

\begin{itemize}
    \item Pada tahun 1972 NIST meminta standard enkripsi \textit{cipher} blok.
    \item Pada tahun 1974 Horst Feistel di IBM mendesain cipher blok \textit{Lucifer} (panjang kunci 128 bit, Panjang blok 128 bit)
    \item \textit{Lucifer} kemudian berkembang menjadi DES dengan beberapa masukan dari NSA (\textit{National Security Agency}):
    \begin{itemize}
        \item panjang kunci 56 bit, ukuran blok 64 bit
        \item 16 putaran.
    \end{itemize}
    \item Tahun 1976 DES disetujui oleh \textit{National Bureau of Standard} (NBS) setelah penilaian kekuatannya oleh NSA.
    \item Sejak itu DES diimplementasikan secara luas, baik \textit{software} maupun \textit{hardware}.
\end{itemize}

\end{document}
